
\makeatletter%
\newcommand\controllaClasse[2]{%
    \@ifclassloaded{article}{#1}{#2}%
}%
\makeatother%

% \controllaClasse{% Classe «article»
% }{% Classe «beamer»
% }%

\controllaClasse{% Classe «article»
    \usepackage[%
        % showframe,%
        % showcrop,%
        hmargin=1cm,%
        bottom=2cm,%
        top=2cm,%
        twocolumn%
    ]{geometry}%
    % \usepackage{multicol} % Pacchetto per gestire le colonne, specie quando si ha un documento misto
    \setlength{\columnsep}{0.75cm}%
}{% Classe «beamer»
    % \usepackage{showframe}
    \setlength\parskip{7.5pt} % Impone la distanza tra due capoversi uguale a «7.5pt» ove il valore di base nella classe Beamer è «0pt»
    % \setlength{\parindent}{15pt} % Impone il rientro al valore di base 15pt, poiché la classe Beamer la pone nulla se non altrimenti specificato
}%

\controllaClasse{% Classe «article»
    \usepackage[british]{babel}
    % \usepackage[utf8]{inputenc} % Se si sta lavorando con LuaLaTeX o XeLaTeX allora si può commentare mentre con pdfLaTeX è necessario    
}{% Classe «beamer»
    \usepackage[italian]{babel}
}%

\usepackage[T1]{fontenc}
% \usepackage[T1]{tipa}
\usepackage{microtype} % Migliora la giustificazione del testo

% Cambia il fonte usato nel documento secondo le filze .ttf salvate in «path»
\usepackage{fontspec} % È necessario il compilatore «LuaLaTeX» o «XeLaTeX»
\controllaClasse{%
    \setmainfont{Vollkorn}[%
        Path = ../Riferimenti/,	        % Percorso
        Extension = .ttf,               % Estensione
        UprightFont = *-Regular,        % Fonte tondo
        ItalicFont = *-Italic,          % Fonte corsivo
        BoldFont = *-Bold,              % Fonte grassetto
        BoldItalicFont = *-BoldItalic   % Fonte grassetto corsivato
    ]%
}{%
    \setsansfont{Vollkorn}[%
        Path = ../Riferimenti/,	        % Percorso
        Extension = .ttf,               % Estensione
        UprightFont = *-Regular,        % Fonte tondo
        ItalicFont = *-Italic,          % Fonte corsivo
        BoldFont = *-Bold,              % Fonte grassetto
        BoldItalicFont = *-BoldItalic   % Fonte grassetto corsivato
    ]%
}%
% Le quattro filze per il fonte sono salvati in 
% «./Percorso/della/Filza/Fonte-Regular.ttf»
% «./Percorso/della/Filza/Fonte-Italic.ttf»
% «./Percorso/della/Filza/Fonte-Bold.ttf»
% «./Percorso/della/Filza/Fonte-BoldItalic.ttf»

\usepackage[backend=biber,style=numeric-comp]{biblatex}%
\addbibresource{../Riferimenti/Bibliografia.bib}
\nocite{*} % Cite all entries inside the «.bib» file
\usepackage{csquotes}

% Pacchetti per vari comandi matematici
\usepackage{amsmath,
            amsfonts,
            amsthm,
            amssymb,
            amsthm}

\usepackage{bbold}     % Per es. l'1 con grassetto da lavagna
\usepackage{mathtools} % Per es. \xrightleftharpoons
\usepackage{empheq}
\usepackage{cancel}

\usepackage{nicematrix,
            array}

\usepackage{graphicx} % Pacchetto per le immagini
\usepackage[shortlabels]{enumitem} % Pacchetto per gestire le liste
\setlist[itemize]{label=$\diamond$}

\usepackage{tikz}
\usetikzlibrary{
    tikzmark,
    positioning,
    arrows.meta,
    fit,
    calc,
    decorations.pathreplacing,
    calligraphy,
    spy,
}
\usepackage{pgfplots}
\usepgfplotslibrary{groupplots}
\pgfplotsset{compat=1.18}

\usepackage{caption,    % Pacchetto per le didascalie dentro le minipagine
            subcaption} % Pacchetto per avere due figure affiancate
\usepackage{float,
            stfloats}

\usepackage{xcolor}   % Pacchetto per colorare il testo
\usepackage{hyperref} % Pacchetto per i collegamenti ipertestuali
\newcommand{\teqref}[2]{\hyperref[#1]{#2}} % Riferimenti applicati a un testo

% Il comando da usare per questo pacchetto è \cref{} anziché \eqref{} o \ref{}
\controllaClasse{% Classe «article»    
    \usepackage{cleveref} % Pacchetto per citare multiple equazioni
    %
    \crefrangeformat{equation}{(#3#1#4--#5#2#6)}%
    \crefmultiformat{equation}{(#2#1#3}{, #2#1#3)}{, #2#1#3}{, #2#1#3)}%
    \crefrangemultiformat{equation}{(#3#1#4--#5#2#6}{, #3#1#4--#5#2#6)}{, #3#1#4--#5#2#6}{, #3#1#4--#5#2#6)}%
    %
    \crefrangeformat{enumi}{#3#1#4--#5#2#6}%
    \crefmultiformat{enumi}{#2#1#3}{ and #2#1#3}{, #2#1#3}{ and #2#1#3}%
    %
    \crefrangeformat{figure}{#3#1#4--#5#2#6}%
    \crefmultiformat{figure}{#2#1#3}{ and #2#1#3}{, #2#1#3}{ and #2#1#3}%
    \crefrangemultiformat{figure}{#3#1#4--#5#2#6}{ and #3#1#4--#5#2#6}{, #3#1#4--#5#2#6}{ and #3#1#4--#5#2#6}%
    %
    \crefrangeformat{section}{#3#1#4--#5#2#6}%
    \crefmultiformat{section}{#2#1#3}{ and #2#1#3}{, #2#1#3}{ and #2#1#3}%
}{% Classe «beamer»
    \usepackage{cleveref} % Pacchetto per citare multiple equazioni
    %
    \crefrangeformat{enumi}{#3#1#4--#5#2#6}%
    \crefmultiformat{enumi}{#2#1#3}{ e #2#1#3}{, #2#1#3}{ e #2#1#3}%
}%

\usepackage{siunitx}
\sisetup{
    inter-unit-product=\ensuremath{{}\cdot{}},
    scientific-notation=false,
}

\controllaClasse{}{% Classe «beamer»
    \beamertemplatenavigationsymbolsempty % Elimina la barra di navigazione in fondo alle diapo
    %
    \usepackage{ragged2e} % Pacchetto per giustificare il testo
    \justifying%
    %
    \renewcommand{\implies}{\DOTSB\Longrightarrow} % Nella definizione standarda «\implies» ha forma «\DOTSB \;\Longrightarrow \;», ove «\DOTSB» corrisponde a «\relax» vale a dire che non fa niente; tuttavia beamer [mannaggia non so perché] frigna che i due «\;» sono invalidi (in generale in modalità matematica par essere cosí dato che non li accetta), quindi è necessario ridefinire «\implies» senza i due simboli cosí da non avere errori imprevisti
    %
    % Settaggio delle didascalie
    % \captionsetup{aboveskip=5pt plus 0pt minus 0pt,           % Spazio sopra la didascalia
    %               belowskip=-2.5pt plus 2.5pt minus 2.5pt}    % Spazio sotto la didascalia
}%